% =====================================================================
%  Kết luận (không đánh số chương)
% =====================================================================
\chapter*{Kết luận}
\addcontentsline{toc}{chapter}{Kết luận}
\markboth{Kết luận}{Kết luận}

\section*{Kết quả đạt được}
Tiểu luận đã đi hết chuỗi kỹ thuật của một hệ thống thông minh: lược sử AI, học máy trên dữ liệu bảng và văn bản, CNN trên ảnh, RNN trên chuỗi thời gian, đồ thị tri thức Neo4j với RAG và tác tử
chạy LLM cục bộ, hệ đa tác tử, mạng nơ-ron đồ thị, học tăng cường, và một Web App ghép các phần lại. Bảng~\ref{tab:kl-best} lấy trực tiếp từ tệp kết quả của từng chương, nêu mô hình tốt nhất theo chỉ số chính trên mỗi bộ dữ liệu.

\bang{ch10_best}{Giá trị tốt nhất trên tập kiểm thử của từng bộ dữ liệu, chỉ để đối chiếu; mô hình đại diện của mỗi chương được chọn trên tập kiểm định.}{kl-best}

Ở Chương 8, tác tử dựa trên mô hình đạt giá trị \SL{ch8_q10/tham_chieu/model/exact_mean} mỗi bước ở cấu hình tham chiếu và \SL{ch8_q10/giam_5/model/exact_mean} khi $\delta=5\%$, bằng mức tối ưu ở cả năm kịch bản, trong khi Q-learning gốc đạt
\SL{ch8_q10/giam_5/q/n_opt}/10 seed tối ưu ở $\delta=5\%$.

\section*{Bài học}
\begin{itemize}
  \item \textbf{Đường cơ sở quyết định cách đọc kết quả.} Một con số chỉ có nghĩa khi đặt cạnh mô hình đơn giản nhất hoặc mức tối ưu tính được. Ở Chương 8, nếu thiếu mốc tối ưu \SL{ch8_model/d20/optimal}, không thể biết tác tử luật đã tối ưu sẵn và RL chỉ có thể hòa.
  \item \textbf{Đánh giá phải tách khỏi huấn luyện.} Chỉ chọn mô hình trên tập kiểm định, dùng tập kiểm thử một lần, và lặp nhiều seed khi kết quả phụ thuộc vào ngẫu nhiên.
  \item \textbf{Chỉ số phải khớp mục tiêu.} Huấn luyện theo doanh thu làm lợi nhuận tụt còn \SL{ch8_g/doanh_thu/eval_mean}. Hệ thống tối ưu đúng cái được đo, không phải cái người thiết kế muốn.
  \item \textbf{Giới hạn phần cứng là một phần của thiết kế.} Với GPU 8~GB, việc chọn kích thước mô hình, cách xếp hàng GPU và mức dùng bộ nhớ quyết định việc gì chạy được.
\end{itemize}

\section*{Hạn chế}
Mô hình được huấn luyện trên đủ dữ liệu, nhưng phần đánh giá LLM chỉ dùng 50 câu mỗi miền ở Chương 5 và Chương 6 nên sai số mẫu còn lớn; nhiều thực nghiệm chỉ chạy một seed với số epoch vừa đủ, nên chưa nêu được độ lệch giữa các lần chạy. LLM cục bộ (mặc định Qwen2.5-7B qua Ollama) nhỏ hơn nhiều so với mô hình trên đám mây nên kết quả ở Chương 5 và 6 chỉ phản ánh cỡ mô hình này. Dữ liệu ở Chương 8 và ca lâm sàng ở Chương 6 là dữ liệu tổng hợp. Các kết quả
về y tế và đầu tư chỉ mang tính học thuật.

\section*{Hướng phát triển}
Có ba hướng rõ ràng: thêm tồn kho và nhiều mức giá cho mô hình thị trường, đối thủ có trí nhớ dài hơn; đưa quan sát từ đồ thị tri thức vào tác tử để đánh giá trên câu hỏi đa bước dài hơn; và theo dõi mô hình sau triển khai trên Web App để phát hiện suy giảm chất lượng.
